\chapter{Annexes}

\subsubsection*{Grille d'entretien \hrulefill}

        \subsection*{Question directrice : }

        \guillemetleft J'aimerais que vous me parliez de ce que vous piratez, depuis quand c'est un mode de consommation culturelle pour vous, que vous me racontiez comment vous le vivez au quotidien, si vous le partagez avec d'autres personnes et comment ça se traduit dans les différents domaines de la vie, professionnelle, personnelle, hobbies, dans vos études...\guillemetright

        \subsection*{Sujets que j'aimerais aborder : }
        
         \begin{itemize}

         \subsubsection*{Histoire personnelle :}
         
             \item Vous rappelez-vous de quand et de comment avez vous découvert le piratage ? Que recherchiez vous, que piratiez vous ?

            \item Saviez vous à l'époque que c'était du piratage ? Comment avez vous découvert que c'était du piratage ?

            \item Aujourd'hui, quelle est la fréquence à laquelle vous piratez du contenu culturel ?

            \item Vous rappelez-vous des périodes de votre vie où vous avez plus piraté ?
             
            \item Avez-vous aidé des proches moins à l'aise avec la technologie à accéder à des contenus piratés ? Avez-vous été dans une démarche de partage de cette pratique ?

        \subsubsection*{Questions sur les émotions :}
        
            \item Quelles sont les émotions que vous procure le fait d'avoir accès à un contenu culturel par le piratage ?

            \item Avez-vous déjà parlé de piratage à votre famille ? A vos ami.e.s ?

            \item Avez vous revendiqué une identité de pirate ou de hacker ?

            \item Qu'est ce qui vous a pousser à commencer à pirater du contenu culturel.
            
            \item Piratez vous pour les mêmes raisons qu'au début. 

        \subsubsection*{Questions techniques :}
        
             \item Avec quels outils piratiez vous (simple téléchargement, sites spécialisés, forums, discord, Reddit, torrenting, IPTV,  applications spécialisées)

             \item Est-ce que vous stockez les médias que vous piratez ? Comment stockez vous les médias que vous piratez ?

             \item Est ce que votre dépôt de contenu piraté suit une organisation particulière, ou un classement ?

             \item Est ce que votre dépôt de contenu piraté est accessible à distance, à vous ou à vos proches ?
             
         \end{itemize}

         pratique plus libre dans les années 2010 ?

\newpage

\subsubsection*{Pacte d'entretien d'information et de recueil de consentement}
    \hrulefill

    \subsection*{Question directrice}

    \guillemetleft J'aimerais que vous me parliez de ce que vous piratez, depuis quand c'est un mode de consommation culturelle pour vous, que vous me racontiez comment vous le vivez au quotidien, si vous le partagez avec d'autres personnes et comment ça se traduit dans les différents domaines de la vie, professionnelle, personnelle, dans vos hobbies, dans vos études...\guillemetright

    \subsection*{Collecte de données personnelles}

    Pour pouvoir contextualiser notre entretien, je vous interrogerai sur des sujets comme la catégorie socioprofessionnelle de vos parents ainsi que leur niveau d'études, votre catégorie socioprofessionnelle et votre parcours scolaire, votre utilisation des ordinateurs, des smartphones, et de votre niveau de maîtrise en informatique.\\

    Un enregistrement audio est également nécessaire à la production d'une transcription qui sera le matériel cité dans mon travail.

    \subsection*{Publication, modification et restitution}

    La transcription et les extraits utilisés seront publiés sous le nom que vous souhaitez, sous pseudonyme, ou avec une anonymisation totale.\\

    Vous pouvez également à tout moment rajouter, corriger, modifier ou supprimer un élément de l'entretien sur lequel vous souhaiteriez ajouter une précision, ou exprimer un changement d'avis.\\

    Comme ce travail se nourrit de vos expériences et de votre vécu, il est normal de vous laisser le choix d'avoir accès soit à une restitution des conclusions du travail final, soit au manuscrit du mémoire si cela vous intéresse.\\
    \vspace{0.5cm}
    
    Pour toutes questions ou modifications, vous pouvez me contacter à cette adresse :
    \vspace{0.1cm}
    
    jules.musquin@chartes.psl.eu